% Options for packages loaded elsewhere
\PassOptionsToPackage{unicode}{hyperref}
\PassOptionsToPackage{hyphens}{url}
%
\documentclass[
  10pt,
  letterpaper,
]{article}
\usepackage{amsmath,amssymb}
\usepackage{iftex}
\ifPDFTeX
  \usepackage[T1]{fontenc}
  \usepackage[utf8]{inputenc}
  \usepackage{textcomp} % provide euro and other symbols
\else % if luatex or xetex
  \usepackage{unicode-math} % this also loads fontspec
  \defaultfontfeatures{Scale=MatchLowercase}
  \defaultfontfeatures[\rmfamily]{Ligatures=TeX,Scale=1}
\fi
\usepackage{lmodern}
\ifPDFTeX\else
  % xetex/luatex font selection
\fi
% Use upquote if available, for straight quotes in verbatim environments
\IfFileExists{upquote.sty}{\usepackage{upquote}}{}
\IfFileExists{microtype.sty}{% use microtype if available
  \usepackage[]{microtype}
  \UseMicrotypeSet[protrusion]{basicmath} % disable protrusion for tt fonts
}{}
\makeatletter
\@ifundefined{KOMAClassName}{% if non-KOMA class
  \IfFileExists{parskip.sty}{%
    \usepackage{parskip}
  }{% else
    \setlength{\parindent}{0pt}
    \setlength{\parskip}{6pt plus 2pt minus 1pt}}
}{% if KOMA class
  \KOMAoptions{parskip=half}}
\makeatother
\usepackage{xcolor}
\usepackage[margin=1in]{geometry}
\usepackage{color}
\usepackage{fancyvrb}
\newcommand{\VerbBar}{|}
\newcommand{\VERB}{\Verb[commandchars=\\\{\}]}
\DefineVerbatimEnvironment{Highlighting}{Verbatim}{commandchars=\\\{\}}
% Add ',fontsize=\small' for more characters per line
\newenvironment{Shaded}{}{}
\newcommand{\AlertTok}[1]{\textcolor[rgb]{1.00,0.00,0.00}{\textbf{#1}}}
\newcommand{\AnnotationTok}[1]{\textcolor[rgb]{0.38,0.63,0.69}{\textbf{\textit{#1}}}}
\newcommand{\AttributeTok}[1]{\textcolor[rgb]{0.49,0.56,0.16}{#1}}
\newcommand{\BaseNTok}[1]{\textcolor[rgb]{0.25,0.63,0.44}{#1}}
\newcommand{\BuiltInTok}[1]{\textcolor[rgb]{0.00,0.50,0.00}{#1}}
\newcommand{\CharTok}[1]{\textcolor[rgb]{0.25,0.44,0.63}{#1}}
\newcommand{\CommentTok}[1]{\textcolor[rgb]{0.38,0.63,0.69}{\textit{#1}}}
\newcommand{\CommentVarTok}[1]{\textcolor[rgb]{0.38,0.63,0.69}{\textbf{\textit{#1}}}}
\newcommand{\ConstantTok}[1]{\textcolor[rgb]{0.53,0.00,0.00}{#1}}
\newcommand{\ControlFlowTok}[1]{\textcolor[rgb]{0.00,0.44,0.13}{\textbf{#1}}}
\newcommand{\DataTypeTok}[1]{\textcolor[rgb]{0.56,0.13,0.00}{#1}}
\newcommand{\DecValTok}[1]{\textcolor[rgb]{0.25,0.63,0.44}{#1}}
\newcommand{\DocumentationTok}[1]{\textcolor[rgb]{0.73,0.13,0.13}{\textit{#1}}}
\newcommand{\ErrorTok}[1]{\textcolor[rgb]{1.00,0.00,0.00}{\textbf{#1}}}
\newcommand{\ExtensionTok}[1]{#1}
\newcommand{\FloatTok}[1]{\textcolor[rgb]{0.25,0.63,0.44}{#1}}
\newcommand{\FunctionTok}[1]{\textcolor[rgb]{0.02,0.16,0.49}{#1}}
\newcommand{\ImportTok}[1]{\textcolor[rgb]{0.00,0.50,0.00}{\textbf{#1}}}
\newcommand{\InformationTok}[1]{\textcolor[rgb]{0.38,0.63,0.69}{\textbf{\textit{#1}}}}
\newcommand{\KeywordTok}[1]{\textcolor[rgb]{0.00,0.44,0.13}{\textbf{#1}}}
\newcommand{\NormalTok}[1]{#1}
\newcommand{\OperatorTok}[1]{\textcolor[rgb]{0.40,0.40,0.40}{#1}}
\newcommand{\OtherTok}[1]{\textcolor[rgb]{0.00,0.44,0.13}{#1}}
\newcommand{\PreprocessorTok}[1]{\textcolor[rgb]{0.74,0.48,0.00}{#1}}
\newcommand{\RegionMarkerTok}[1]{#1}
\newcommand{\SpecialCharTok}[1]{\textcolor[rgb]{0.25,0.44,0.63}{#1}}
\newcommand{\SpecialStringTok}[1]{\textcolor[rgb]{0.73,0.40,0.53}{#1}}
\newcommand{\StringTok}[1]{\textcolor[rgb]{0.25,0.44,0.63}{#1}}
\newcommand{\VariableTok}[1]{\textcolor[rgb]{0.10,0.09,0.49}{#1}}
\newcommand{\VerbatimStringTok}[1]{\textcolor[rgb]{0.25,0.44,0.63}{#1}}
\newcommand{\WarningTok}[1]{\textcolor[rgb]{0.38,0.63,0.69}{\textbf{\textit{#1}}}}
\setlength{\emergencystretch}{3em} % prevent overfull lines
\providecommand{\tightlist}{%
  \setlength{\itemsep}{0pt}\setlength{\parskip}{0pt}}
\setcounter{secnumdepth}{-\maxdimen} % remove section numbering
\usepackage{xurl}
\setlength{\emergencystretch}{3em}
\usepackage{microtype}
\ifLuaTeX
  \usepackage{selnolig}  % disable illegal ligatures
\fi
\usepackage{bookmark}
\IfFileExists{xurl.sty}{\usepackage{xurl}}{} % add URL line breaks if available
\urlstyle{same}
\hypersetup{
  hidelinks,
  pdfcreator={LaTeX via pandoc}}

\author{}
\date{}

\begin{document}

\section{BA2-T2 - Semantic operator, participation-role and
scoped-modifier vocabulary pressure
test}\label{ba2-t2---semantic-operator-participation-role-and-scoped-modifier-vocabulary-pressure-test}

\textbf{Revision:} R1

\textbf{Status:} COMPLETED / PROVISIONAL PASS WITH
VOCABULARY-ARCHITECTURE REFINEMENT / BA2 NOT CLOSED

\textbf{Repository baseline reviewed:}
\texttt{f87d05e5ea1cee969246e5eae1dd73b8b6c3a5a1}

\textbf{Phase:} BA2 - Relations and canonical action vocabulary

\textbf{BA0:} CLOSED

\textbf{BA1:} CLOSED

\textbf{BA3:} NOT STARTED

\subsection{1. Question and test
discipline}\label{question-and-test-discipline}

BA2-T2 asks only:

\begin{quote}
What is the smallest reusable methodology-neutral vocabulary
architecture needed to instantiate the BA2-T1 proposition shape across
the current governed corpora without copying authoring verbs, importing
method-specific taxonomies or hiding reusable semantics inside generic
modifiers?
\end{quote}

The starting structural candidate is:

\begin{Shaded}
\begin{Highlighting}[]
\NormalTok{BAProposition}
\NormalTok{|{-} semanticOperator   1}
\NormalTok{|{-} participation      1..*}
\NormalTok{|    |{-} role          1}
\NormalTok{|    \textasciigrave{}{-} term          1}
\NormalTok{\textasciigrave{}{-} scopedModifier     0..*}
\end{Highlighting}
\end{Shaded}

BA2-T2 does not close BA2, design BA3 provenance, define views, create a
formal STRIDE overlay, define Common Finding, or create ThreatForge
implementation classes.

The pressure discipline is:

\begin{Shaded}
\begin{Highlighting}[]
\NormalTok{recurring lexical verb}
\NormalTok{        !=}
\NormalTok{canonical semantic operator}

\NormalTok{reusable role label}
\NormalTok{        !=}
\NormalTok{globally valid role in every proposition}

\NormalTok{local qualifier}
\NormalTok{        !=}
\NormalTok{license for an untyped modifier bag}
\end{Highlighting}
\end{Shaded}

\subsection{2. Evidence basis}\label{evidence-basis}

The test reuses four already-governed evidence families without changing
them:

\begin{enumerate}
\def\labelenumi{\arabic{enumi}.}
\tightlist
\item
  the closed BA1 identity ontology
  (\texttt{BAReferent\ +\ BAProposition});
\item
  the BA2-T1 n-ary role-bound proposition candidate;
\item
  facial-access FR-3.4 plus D-3.3/D-3.4/D-3.5 and specialized
  constraints;
\item
  the independent order-fulfillment corpus, including reservation,
  payment, fulfillment, external providers, lifecycle and physical
  handoff reuse.
\end{enumerate}

The BA0-T2 STRIDE-oriented transfer projection is used only as a bounded
consumer probe. It cannot donate STRIDE/DFD vocabulary to Base Analysis.

\subsection{3. Pressure A - copy document predicates or normalize
semantic
operators?}\label{pressure-a---copy-document-predicates-or-normalize-semantic-operators}

The order corpus contains surface predicates such as \texttt{creates},
\texttt{records}, \texttt{increases}, \texttt{updates}, \texttt{reads},
\texttt{produces}, \texttt{releases}, \texttt{decreasesReserved},
\texttt{sends}, \texttt{communicatesWith}, \texttt{maps},
\texttt{queries}, \texttt{requestsCapture}, \texttt{requestsVoid},
\texttt{confirmsPick}, \texttt{confirmsHandoff} and
\texttt{requestsStockIssue}.

A naive BA vocabulary could simply copy all of them.

That alternative fails for three reasons.

First, documentation predicates are authoring vocabulary, not
automatically canonical Base Analysis semantics. The documentation layer
itself kept the FunctionalPredicate vocabulary extensible and deferred
subject/object binding until Base Analysis work.

Second, some surface verbs express the same broader project-semantic
relation while other similar words conceal materially different
semantics. \texttt{requestsCapture} and \texttt{requestsVoid}, for
example, are both requests but differ in governed financial meaning;
\texttt{confirmsHandoff} is not interchangeable with carrier acceptance.

Third, a method consumer needs stable semantics despite lexical change
or provider replacement.

\subsubsection{Disposition}\label{disposition}

\begin{Shaded}
\begin{Highlighting}[]
\NormalTok{Copy source predicate list as BA canonical vocabulary: REJECTED}
\NormalTok{Stable semantic operator key with normative meaning: REQUIRED}
\NormalTok{Lexical label/synonym as semantic authority: REJECTED}
\end{Highlighting}
\end{Shaded}

\subsection{4. Pressure B - does an operator-family layer
help?}\label{pressure-b---does-an-operator-family-layer-help}

The corpora repeatedly distinguish behavioral assertions, stable
relationships, constraints and reusable classification.

A provisional family facet (\texttt{ACTION}, \texttt{RELATION},
\texttt{CONSTRAINT}, \texttt{CLASSIFICATION}) helps organize the
registry and narrow compatible role/modifier contracts.

However, the family alone is too coarse. A consumer cannot treat
\texttt{transfer}, \texttt{transition} and \texttt{produce} as
semantically interchangeable merely because all are actions. Likewise
\texttt{correlate}, \texttt{realize} and \texttt{consumeService} are not
interchangeable relations.

\subsubsection{Disposition}\label{disposition-1}

\begin{Shaded}
\begin{Highlighting}[]
\NormalTok{Exact operator key: REQUIRED}
\NormalTok{Operator{-}family facet: SUPPORTED AS PROVISIONAL ORGANIZING FACET}
\NormalTok{Operator family as sufficient assertion semantics: REJECTED}
\NormalTok{Family facet as new BAE identity family: REJECTED}
\end{Highlighting}
\end{Shaded}

The family facet remains falsifiable and may collapse later if
regression shows it adds no semantic value.

\subsection{5. Pressure C - semantic distinctions that must survive
normalization}\label{pressure-c---semantic-distinctions-that-must-survive-normalization}

The pressure test intentionally tries to merge several recurring
relations.

\subsubsection{5.1 Correlation versus
reference}\label{correlation-versus-reference}

\texttt{Reservation} references one \texttt{OrderEvaluation};
payment/reservation/provider results are correlated to the current
evaluation; facial capture is correlated with the correct
\texttt{RecognitionRequest}.

A generic \texttt{relatedTo} operator would erase whether identity
matching is required for acceptance versus whether one object merely
references another.

\begin{Shaded}
\begin{Highlighting}[]
\NormalTok{correlation != generic reference}
\end{Highlighting}
\end{Shaded}

\textbf{Disposition:} distinction
\texttt{SUPPORTED\ /\ SEED\ KEYS\ CANDIDATE}.

\subsubsection{5.2 Service consumption versus
ownership}\label{service-consumption-versus-ownership}

Facial access consumes transport while the project does not own/manage
the underlying transport. Order fulfillment consumes cross-MR
capabilities without gaining a second owner.

\begin{Shaded}
\begin{Highlighting}[]
\NormalTok{consumeService != ownOrManage}
\end{Highlighting}
\end{Shaded}

Collapsing them would directly lose the responsibility boundary that the
corpora deliberately govern.

\textbf{Disposition:} distinction
\texttt{FORCED\ BY\ CURRENT\ EVIDENCE}.

\subsubsection{5.3 Request acceptance versus physical
handoff}\label{request-acceptance-versus-physical-handoff}

The order corpus explicitly rejects treating carrier request acceptance
as physical handoff. Stock issue and payment capture depend on the
governed physical milestone.

A generic \texttt{completed} or \texttt{accepted} relation would erase a
material business-event distinction.

\textbf{Disposition:} distinct referent/proposition semantics
\texttt{FORCED\ BY\ CURRENT\ EVIDENCE}.

\subsubsection{5.4 Realization versus
qualification}\label{realization-versus-qualification}

Ethernet/Wi-Fi can change while the same abstract connectivity meaning
persists. Treating realization as a free-text modifier would weaken
selective change impact and reuse.

\textbf{Disposition:} \texttt{realize} relation remains explicit
\texttt{CANDIDATE}; generic modifier representation \texttt{REJECTED}
for reusable realizations.

\subsection{6. Pressure D - seed operator
set}\label{pressure-d---seed-operator-set}

BA2-T2 then asks whether a small seed set can cover the tested evidence
without pretending to be exhaustive.

The surviving candidate keys are grouped as follows:

\begin{Shaded}
\begin{Highlighting}[]
\NormalTok{ACTION}
\NormalTok{  transfer}
\NormalTok{  produce}
\NormalTok{  create}
\NormalTok{  observe}
\NormalTok{  transition}

\NormalTok{RELATION}
\NormalTok{  correlate}
\NormalTok{  reference}
\NormalTok{  dependOn}
\NormalTok{  consumeService}
\NormalTok{  realize}
\NormalTok{  assignResponsibility}
\NormalTok{  ownOrManage}

\NormalTok{CONSTRAINT}
\NormalTok{  constrain}

\NormalTok{CLASSIFICATION}
\NormalTok{  classify}
\end{Highlighting}
\end{Shaded}

These names are provisional readable keys. Their exact machine key
format and lexical aliases remain open.

\subsubsection{Why this is not a verb taxonomy
closure}\label{why-this-is-not-a-verb-taxonomy-closure}

The seed set is a regression target. A key survives only if collapsing
it would lose project semantics required by a corpus or consumer. New
document verbs do not automatically extend it, and later evidence may
split or merge keys.

\subsubsection{Disposition}\label{disposition-2}

\begin{Shaded}
\begin{Highlighting}[]
\NormalTok{Seed operator set: PROVISIONAL CANDIDATE}
\NormalTok{Exhaustive operator vocabulary: NOT CLAIMED}
\NormalTok{BA2 closure: NOT REACHED}
\end{Highlighting}
\end{Shaded}

\subsection{7. Pressure E - global roles or operator-scoped role
contracts?}\label{pressure-e---global-roles-or-operator-scoped-role-contracts}

A global role list looks attractive because names such as
\texttt{source}, \texttt{destination}, \texttt{result} and
\texttt{context} recur.

But global validity is unsafe. \texttt{source} in a transfer is not BA3
source provenance. \texttt{target} could mean destination, constraint
target or analysis-method target. \texttt{context} can mean correlation
identity, state observation context or applicability context.

The better lower bound is:

\begin{Shaded}
\begin{Highlighting}[]
\NormalTok{reusable role concept/key}
\NormalTok{        +}
\NormalTok{operator{-}scoped compatibility/cardinality contract}
\end{Highlighting}
\end{Shaded}

Example contracts:

\begin{Shaded}
\begin{Highlighting}[]
\NormalTok{transfer:}
\NormalTok{  source       1}
\NormalTok{  destination  1}
\NormalTok{  content      1..*}
\NormalTok{  request/context 0..*}

\NormalTok{transition:}
\NormalTok{  subject      1}
\NormalTok{  fromState    0..1}
\NormalTok{  toState      1}

\NormalTok{consumeService:}
\NormalTok{  consumer     1..*}
\NormalTok{  service      1}
\NormalTok{  provider     0..1 (or governed cardinality later)}

\NormalTok{classify:}
\NormalTok{  classifiedReferent 1}
\NormalTok{  semanticKind       1..*}
\end{Highlighting}
\end{Shaded}

\subsubsection{Disposition}\label{disposition-3}

\begin{Shaded}
\begin{Highlighting}[]
\NormalTok{One context{-}free global role contract: REJECTED}
\NormalTok{Reusable role keys/concepts: SUPPORTED}
\NormalTok{Operator{-}scoped compatibility/cardinality: REQUIRED}
\NormalTok{Complete role matrix: OPEN / REGRESSION TARGET}
\end{Highlighting}
\end{Shaded}

\subsection{8. Pressure F - participant or
modifier?}\label{pressure-f---participant-or-modifier}

BA2-T1 left the exact modifier boundary open. BA2-T2 tests it against
condition, state, failure, applicability, realization, atomicity,
concurrency and polarity.

\subsubsection{8.1 Modifier admissibility
criterion}\label{modifier-admissibility-criterion}

An embedded modifier is acceptable only if it is local to one
proposition/binding, introduces no independent participant set, needs no
independent assertion provenance/review identity, and is not reused as
project meaning elsewhere.

If those conditions are violated, the semantics must be represented as a
referent/proposition rather than hidden inside the modifier.

\subsubsection{8.2 Semantics that fit local
modifiers}\label{semantics-that-fit-local-modifiers}

Current evidence supports candidate modifier kinds for:

\begin{itemize}
\tightlist
\item
  polarity;
\item
  local condition/applicability;
\item
  quantification/cardinality;
\item
  temporal scope;
\item
  local completion/failure outcome;
\item
  atomicity;
\item
  concurrency;
\item
  idempotency.
\end{itemize}

\subsubsection{8.3 Semantics that should not default to
modifiers}\label{semantics-that-should-not-default-to-modifiers}

Correlation, realization, responsibility/authority, service consumption
and reusable constraints all have reusable participants or independent
assertion pressure in the reviewed corpora.

They should therefore remain explicit propositions/participations when
that semantic pressure exists.

\subsubsection{Disposition}\label{disposition-4}

\begin{Shaded}
\begin{Highlighting}[]
\NormalTok{Untyped/free{-}text modifier bag: REJECTED}
\NormalTok{Normalized modifier kinds: SUPPORTED / CANDIDATE}
\NormalTok{Modifier promotion rule: REQUIRED}
\NormalTok{Exact modifier encoding: OPEN}
\end{Highlighting}
\end{Shaded}

\subsection{9. Pressure G - explicit
polarity}\label{pressure-g---explicit-polarity}

The corpora contain important negative semantics:

\begin{itemize}
\tightlist
\item
  project does not own/manage underlying transport;
\item
  incomplete delivery is not success;
\item
  a concurrent request must not reserve already frozen quantity;
\item
  raw provider/carrier states must not directly become governed project
  state;
\item
  an indeterminate/rejected carrier request alone must not count as
  physical handoff.
\end{itemize}

Creating \texttt{notOwn}, \texttt{notAccept}, \texttt{notHandoff},
\texttt{notReserve} keys would make operator vocabulary explode around
linguistic negation.

A normalized polarity field/modifier separates assertion polarity from
operator identity.

\subsubsection{Disposition}\label{disposition-5}

\begin{Shaded}
\begin{Highlighting}[]
\NormalTok{Separate notX operator keys by default: REJECTED}
\NormalTok{Explicit polarity: STRONG CANDIDATE}
\end{Highlighting}
\end{Shaded}

Some normative prohibitions will still be represented as positive
\texttt{constrain} propositions with negative/forbidden outcomes;
polarity does not eliminate constraint semantics.

\subsection{10. Pressure H - reusable referent
classification}\label{pressure-h---reusable-referent-classification}

BA2-T1 established that referent classification cannot be reconstructed
only from contextual roles.

BA2-T2 pressure-tests two alternatives:

\begin{enumerate}
\def\labelenumi{\arabic{enumi}.}
\tightlist
\item
  opaque intrinsic \texttt{type} field on BAReferent;
\item
  classification through a normal BAProposition using a stable
  \texttt{classify} operator.
\end{enumerate}

The proposition form is stronger for the current methodology because
classification itself may be grounded, derived, disputed or changed
while the referent identity remains stable. BA3 can then provenance the
classification assertion without adding a new BA1 type.

Conceptual example:

\begin{Shaded}
\begin{Highlighting}[]
\NormalTok{operator: classify}
\NormalTok{participants:}
\NormalTok{  classifiedReferent {-}\textgreater{} RecognitionCapture}
\NormalTok{  semanticKind       {-}\textgreater{} information{-}resource}
\end{Highlighting}
\end{Shaded}

\subsubsection{Disposition}\label{disposition-6}

\begin{Shaded}
\begin{Highlighting}[]
\NormalTok{Classification inferred only from roles: REJECTED}
\NormalTok{Opaque authoritative referent type field: NOT REQUIRED}
\NormalTok{Classification{-}as{-}BAProposition: SUPPORTED / STRONG CANDIDATE}
\NormalTok{Exact semantic{-}kind taxonomy: OPEN}
\end{Highlighting}
\end{Shaded}

\subsection{11. Pressure I - bounded method
consumer}\label{pressure-i---bounded-method-consumer}

The BA0-T2 STRIDE-oriented consumer needs behavior, source, destination,
information, correlation, responsibility/externality, realization,
failure and project constraints.

The candidate can supply these without importing DFD/STRIDE roots:

\begin{itemize}
\tightlist
\item
  \texttt{transfer} + source/destination/content/context roles;
\item
  \texttt{correlate} where correlation is independently queried;
\item
  \texttt{consumeService}, \texttt{assignResponsibility} and negative
  \texttt{ownOrManage} for responsibility/externality;
\item
  \texttt{realize} for current Ethernet realization;
\item
  local completion/failure modifier or explicit \texttt{constrain}
  proposition as warranted;
\item
  \texttt{classify} for reusable behavior/information/capability
  semantic kind.
\end{itemize}

STRIDE then adds Spoofing/Tampering/Repudiation/Information
Disclosure/DoS/Elevation semantics outside Base Analysis.

\subsubsection{Disposition}\label{disposition-7}

\begin{Shaded}
\begin{Highlighting}[]
\NormalTok{Bounded method{-}consumer mapping: PASS AT BA2{-}T2 CANDIDATE LEVEL}
\NormalTok{STRIDE/DFD vocabulary leakage: NOT REQUIRED}
\NormalTok{Formal method projection contract: DEFERRED TO BA4 / LATER OVERLAY WORK}
\end{Highlighting}
\end{Shaded}

\subsection{12. Pressure J - does BA2-T2 reopen
BA1?}\label{pressure-j---does-ba2-t2-reopen-ba1}

No tested operator, role or modifier requires independent first-class
identity beyond the closed \texttt{BAReferent} and
\texttt{BAProposition} families.

Operator keys, role keys, modifier kinds and controlled semantic-kind
values are vocabulary identities/contracts, not project-semantic BAE
identity families.

If a future case requires one of them to have independent
project-semantic identity plus subtype-specific invariants not
representable over BAReferent/BAProposition, that would trigger the BA1
reopen criterion rather than being silently introduced here.

\subsubsection{Disposition}\label{disposition-8}

\begin{Shaded}
\begin{Highlighting}[]
\NormalTok{Third BA1 identity family: NOT FORCED}
\NormalTok{BA1: REMAINS CLOSED}
\end{Highlighting}
\end{Shaded}

\subsection{13. Candidate vocabulary
architecture}\label{candidate-vocabulary-architecture}

BA2-T2 therefore retains the BA2-T1 structure and refines its vocabulary
contract:

\begin{Shaded}
\begin{Highlighting}[]
\NormalTok{BAProposition}
\NormalTok{|{-} semanticOperatorKey   1}
\NormalTok{|    \textasciigrave{}{-} registry contract:}
\NormalTok{|         normative meaning}
\NormalTok{|         provisional family facet}
\NormalTok{|         compatible roles/cardinalities}
\NormalTok{|         modifier compatibility}
\NormalTok{|{-} participation         1..*}
\NormalTok{|    |{-} roleKey          1}
\NormalTok{|    \textasciigrave{}{-} term             1}
\NormalTok{|{-} polarity              1       [STRONG CANDIDATE]}
\NormalTok{\textasciigrave{}{-} scopedModifier        0..*    [normalized local semantics only]}
\end{Highlighting}
\end{Shaded}

Classification is provisionally represented through
\texttt{BAProposition(classify,\ ...)} rather than by adding a
first-class BAE type.

\subsection{14. Alternatives rejected or
deferred}\label{alternatives-rejected-or-deferred}

\begin{itemize}
\tightlist
\item
  document predicate list as BA vocabulary: \textbf{REJECTED};
\item
  lexical label as semantic identity: \textbf{REJECTED};
\item
  fixed exhaustive operator list: \textbf{NOT CLAIMED / DEFERRED};
\item
  family facet as sufficient semantics: \textbf{REJECTED};
\item
  global role validity without operator contract: \textbf{REJECTED};
\item
  complete compatibility/cardinality matrix: \textbf{OPEN};
\item
  untyped modifier bag: \textbf{REJECTED};
\item
  correlation/realization/responsibility hidden as modifiers by default:
  \textbf{REJECTED};
\item
  negative operator proliferation: \textbf{REJECTED};
\item
  explicit polarity: \textbf{STRONG CANDIDATE};
\item
  classification inferred from roles only: \textbf{REJECTED};
\item
  classification proposition: \textbf{STRONG CANDIDATE};
\item
  exact semantic-kind taxonomy: \textbf{OPEN};
\item
  BA3 provenance fields: \textbf{DEFERRED TO BA3};
\item
  STRIDE/DFD semantics in shared vocabulary: \textbf{REJECTED}.
\end{itemize}

\subsection{15. Falsification
assessment}\label{falsification-assessment}

The candidate would fail if it forced consumers to reinterpret raw prose
to distinguish correlation, responsibility, realization or state
semantics; if operator-scoped roles could not cover both corpora; if
local modifiers routinely needed independent participants/provenance; or
if a bounded method consumer required method-specific shared-core
categories.

No such failure is forced by the current evidence.

However, BA2 is not ready to close because the seed key set and
role/cardinality contracts have not yet been replayed systematically
across both corpora as a regression suite. The lexical-key separation
and modifier-promotion edge cases also remain pressure targets.

\subsection{16. Trial disposition}\label{trial-disposition}

\begin{Shaded}
\begin{Highlighting}[]
\NormalTok{BA2{-}T1 structural candidate          RETAINED}
\NormalTok{Stable semantic operator key         REQUIRED}
\NormalTok{Seed operator set                    PROVISIONAL CANDIDATE}
\NormalTok{Operator family facet                PROVISIONAL / FALSIFIABLE}
\NormalTok{Operator{-}scoped role contracts       REQUIRED}
\NormalTok{Full role/cardinality matrix         OPEN}
\NormalTok{Normalized modifier capability       SUPPORTED / CANDIDATE}
\NormalTok{Explicit polarity                    STRONG CANDIDATE}
\NormalTok{Modifier promotion rule              REQUIRED}
\NormalTok{Classification{-}as{-}proposition        STRONG CANDIDATE}
\NormalTok{Exact semantic{-}kind vocabulary       OPEN}
\NormalTok{Third BA1 identity family            NOT FORCED}
\NormalTok{BA1                                  CLOSED}
\NormalTok{BA2{-}T2                               COMPLETED / PROVISIONAL PASS WITH REFINEMENT}
\NormalTok{BA2                                  STARTED / NOT CLOSED}
\NormalTok{BA3                                  NOT STARTED}
\end{Highlighting}
\end{Shaded}

\subsection{17. Next authorized
microstep}\label{next-authorized-microstep}

Execute only:

\begin{quote}
\textbf{BA2-T3 - cross-corpus regression of the operator/role/modifier
candidate and semantic-key/lexical separation.}
\end{quote}

Do not start BA3, formal STRIDE overlay design, Common Finding schema or
implementation work.

\end{document}
